\documentclass{article}

% Final report completed from FAI_Project_Proposal.tex; original style retained.
\usepackage{PRIMEarxiv}
\usepackage[utf8]{inputenc}
\usepackage[T1]{fontenc}
\usepackage{hyperref}
\usepackage{url}
\usepackage{booktabs}
\usepackage{amsfonts,amsmath,amssymb}
\usepackage{microtype}
\usepackage{fancyhdr}
\usepackage{graphicx}
\usepackage{array}
\hypersetup{hidelinks}
\urlstyle{same}
\graphicspath{{final_assets/}}
% Generated from audited raw data; do not edit.
\newcommand{\CertifiedSolves}{270}
\newcommand{\CreditHalfWR}{\(0.844\)}
\newcommand{\CreditTightQ}{\(10.9\)}
\newcommand{\CreditTightWR}{\(0.703\)}
\newcommand{\CreditTiming}{\(+15.84\;[15.32,\,16.37]\)}
\newcommand{\DoubleTwo}{\(-12.47\;[-13.09,\,-11.83]\)}
\newcommand{\KarmaAuctionGain}{\(+66.96\)}
\newcommand{\KarmaForcedGain}{\(-20.90\)}
\newcommand{\KarmaPopulationReversal}{\(+46.06\;[45.40,\,46.74]\)}
\newcommand{\KarmaUncappedReversal}{\(+29.43\;[28.85,\,30.01]\)}
\newcommand{\MediumFive}{\(-7.17\;[-7.50,\,-6.82]\)}
\newcommand{\MediumTwo}{\(+2.00\;[1.52,\,2.49]\)}
\newcommand{\MildFive}{\(+1.05\;[0.81,\,1.30]\)}
\newcommand{\MixedGreedyJB}{\(0.517\)}
\newcommand{\MixedScoreJB}{\(0.995\)}
\newcommand{\MixedScoreWR}{\(0.897\)}
\newcommand{\PersistentScoreSR}{\(44.5\)}
\newcommand{\RankLedgerEffect}{\(+1.31\;[0.88,\,1.73]\)}
\newcommand{\RankPlanner}{\(+2.58\;[2.19,\,2.98]\)}
\newcommand{\ScorePlanner}{\(+2.24\;[1.95,\,2.52]\)}
\newcommand{\ScoreRuntime}{\(111.5\)}
\newcommand{\TimingHigh}{\(+1.372\;[1.260,\,1.494]\)}
\newcommand{\TimingLow}{\(-0.223\;[-0.350,\,-0.107]\)}
\newcommand{\TimingTrend}{\(+1.595\;[1.454,\,1.750]\)}
\newcommand{\TruthGreedyQ}{\(41.7\)}
\newcommand{\TruthGreedySR}{\(8.44\)}
\newcommand{\TruthScoreSR}{\(7.47\)}
\newcommand{\TruthScoreWR}{\(0.996\)}

\pagestyle{fancy}
\fancyhf{}
\fancyhead[LO]{\small Fairness, efficiency, and strategic behaviour in shared GPU allocation}
\cfoot{\thepage}
\setlength{\headheight}{14pt}
\setlength{\parindent}{0pt}
\setlength{\parskip}{4pt}
\newcommand{\WR}{\mathrm{WR}}
\newcommand{\SR}{\mathrm{SR}_{\Delta}}
\newcommand{\Muni}{M_{\mathrm{uni}}}
\newcommand{\repo}{\url{https://github.com/kalpshah18/Shared-GPU-Allocation}}
\newcommand{\resulttable}[1]{\begin{center}\small\setlength{\tabcolsep}{5pt}\input{final_assets/#1}\end{center}}
\newcommand{\resultfigure}[2]{\begin{center}\includegraphics[width=\linewidth]{#1.pdf}\par\small #2\end{center}}

\title{Fairness, efficiency, and strategic behaviour\\in shared GPU allocation}
\author{Aayush Kuloor \quad Harsh Dhru \quad Raj Modi \quad Kalp Shah \\
  Indian Institute of Technology Gandhinagar}

\begin{document}
\maketitle
\thispagestyle{empty}

\begin{abstract}
We study repeated GPU allocation when users privately know their per-round
values and may misreport them. A seeded simulator compares random allocation,
round-robin, greedy allocation, a history-penalised score rule, and a monetary
Vickrey auction. The main question is whether the score rule's fairness benefit
also limits manipulation. At $\lambda=1$, Score retains truthful welfare ratio
\TruthScoreWR{} under homogeneous demand, but mild exaggeration and selective
reporting can remain profitable. A controlled timing experiment changes only
the focal user's policy. At threshold $\tau=0.75$, increasing $\lambda$ from
$0.5$ to $5$ raises the timing premium by \TimingTrend{} utility over 1,000
rounds (95\% interval); a higher threshold remains costly.
We also evaluate Credit-Vickrey and five history/service extensions:
Score-Cap, Karma, Karma-Cap, Rank and Rank-Cap. Credits do not inherit monetary
truthfulness, and a service bound does not ensure incentive compatibility.
All reported measurements come from saved seed-level outputs. The findings
concern synthetic, unit-demand workloads and specified reporting policies;
they do not establish dynamic strategyproofness or production-cluster performance.
\end{abstract}
\keywords{GPU allocation \and fair resource allocation \and mechanism design
  \and strategic behaviour \and simulation}

\section{Introduction}
Shared GPU clusters must balance useful allocation against fair access.
Selecting the highest reported value can favour users who exaggerate. Serving
everyone equally can waste opportunities when values differ. Allocation
history creates another incentive: winning now can reduce a user's priority
later. This project isolates these effects in a repeated, unit-demand model.

The proposal asked how welfare, fairness and manipulation change with scarcity,
history penalties, heterogeneous values and persistent demand. Our main focus
is the history-penalised Score rule (M4). Greedy's vulnerability and monetary
Vickrey's per-round truthfulness are controls with known theoretical
explanations. M4's cross-round incentives require a more direct investigation.

We ask whether a moderate penalty improves fairness cheaply, which reporting
policies still profit as the penalty grows, and whether withholding low-value
reports preserves useful priority. We then test a renewable credit budget and
service-capped alternatives. Positive gains identify profitable tested
deviations; negative gains only show that the particular policy failed.

\paragraph{Context and scope.}
Tiresias schedules distributed training jobs using attained-service ideas,
while Gavel supports heterogeneous accelerators
\cite{gu2019tiresias,narayanan2020heterogeneity}. Our model omits their job and
placement structure. Work on dynamic resource fairness also shows that adding
allocation history can change incentives \cite{fikioris2021incentives}.
Our aim is a reproducible mechanism comparison, not a new production scheduler
or an unrestricted equilibrium result.

\newpage
\section{Preliminaries and problem formulation}
There are $n$ users, $k<n$ identical GPUs and $T$ rounds, indexed
$t=0,\ldots,T-1$. Each user receives at most one GPU per round. Values
$v_{i,t}$ and reports $\hat v_{i,t}$ lie in $[0,1]$. The allocation
$x_{i,t}\in\{0,1\}$ satisfies $\sum_i x_{i,t}=k$. Monetary payments are
$p_{i,t}\ge0$; non-monetary rules return zero utility payments.
Before round $t$, cumulative service is $a_i(t)=\sum_{s<t}x_{i,s}$ and
consecutive waiting is $q_i(t)$, with both initially zero. After service,
$q_i$ resets to zero; otherwise it increases by one.

\paragraph{Utility and welfare.}
Let $G_i=\sum_t v_{i,t}x_{i,t}$ and $U_i=G_i-\sum_t p_{i,t}$.
Credits are bidding resources, not deductions from $U_i$.
True welfare excludes transfers: $W=\sum_i G_i$.
The unconstrained oracle allocates to the $k$ largest true values each round:
\[
 W^*=\sum_t\sum_{i\in S_t^*}v_{i,t},\qquad \WR=W/W^*.
\]
We measure count fairness using $A_i=a_i(T)$ and benefit fairness using
$B_i=G_i/(T\mu_i)$, where $\mu_i$ is the user's marginal mean value:
\[
 J_A=\frac{(\sum_i A_i)^2}{n\sum_i A_i^2},\qquad
 J_B=\frac{(\sum_i B_i)^2}{n\sum_i B_i^2}.
\]
These are applications of Jain's index \cite{jain1984quantitative}. Equal
counts and equal normalised benefits are different objectives.
We also record Nash gross welfare $\sum_i\log(G_i+10^{-8})$, maximum and
95th-percentile waiting, and starvation
$\SR=(nT)^{-1}\sum_{i,t}\mathbf1\{q_i(t)>\Delta\}$ for
$\Delta=2\lceil n/k\rceil$. Reported waiting metrics use pre-round states;
the formal service guarantee also includes the terminal state.

\paragraph{Paired strategic effects.}
A seed package $\omega$ fixes values, tie-breaking and a strategic set
$S$, with $|S|=\lfloor\rho n\rfloor$. For a focal user $f$,
\[
 M_f=U_f(\sigma_f,\sigma_{-f};\omega)
       -U_f(\mathrm{truth},\sigma_{-f};\omega).
\]
Only the focal policy changes. The statistic $\Muni$ averages up to five
focal users within each seed; confidence intervals resample seeds, not users.
The coalition comparison instead changes all strategic users at once and can
have a different sign. Strategic welfare loss is
$\mathrm{SWL}=(W_{\mathrm{truth}}-W_{\mathrm{strategic}})/W_{\mathrm{truth}}$;
it can be negative.
For truthful mechanism $F$, $\mathrm{PoF}=(W^*-W_F)/W^*$.
We compare these objectives separately rather than using an arbitrary weighted sum.

\section{Methodology}
\subsection{Baseline algorithms and the M4 score}
M1 selects $k$ users uniformly without replacement. M2 serves the next $k$
users in a seeded circular queue. Both ignore reports. M3 selects the largest
reports and is welfare-optimal only when reports are truthful. M4 selects
the largest scores
\[
 s_{i,t}=\frac{\hat v_{i,t}}{(1+a_i(t))^\lambda},\qquad \lambda\ge0.
\]
Thus M4 equals M3 at $\lambda=0$. Each win reduces later scores, creating a
possible reason to suppress low reports. The implementation sorts by score,
then an independent seeded tie key; it never adds jitter to small scores.
After allocation, the runner updates service counts and waits.

M5 selects the largest bids and charges each winner the $(k+1)$-st bid.
With unit demand, quasi-linear monetary utility and no intertemporal budget,
truth is optimal in each round \cite{vickrey1961counterspeculation}:
conditional on other bids, winning yields value minus a fixed threshold price.
This property is not transferred to credit-based or history-dependent rules.

\begin{center}\small\setlength{\tabcolsep}{4pt}
\begin{tabular}{lp{4.8in}}
\toprule Studies & Completed scope \\
\midrule
E0--E3 & Finite validation; scarcity; fairness/strategy frontier; strategic fraction. \\
E4--E8 & Heterogeneity; persistence controls; runtime/memory; inflation; history timing. \\
E9/E10/E12 & Extension settings and capped/uncapped component comparisons. \\
E3b/E3c/E11/E15 & Historical planning checkpoints, superseded by corrected E15 attacks. \\
E13/E18 & Adaptive populations and longer training with held-out policy regret. \\
E14/E16/E17 & Constrained welfare oracle; Rank-history and Karma-accounting diagnostics. \\
E19/E20 & Controlled value-threshold timing and renewable-credit budgets. \\
\bottomrule
\end{tabular}\end{center}

\newpage
\subsection{Reporting policies and a controlled timing test}
The core policies are truth, capped exaggeration
$\min(cv,1)$ for $c\in\{1.25,1.5,2\}$, and always reporting one.
History-triggered policies exaggerate only when the user's allocation count
is below a population threshold. Their population-level comparisons change
opponent policies too, so they do not isolate a timing effect.

E19 adds a simple value threshold $\tau\in\{0.5,0.75,0.9\}$.
\emph{Sandbag} reports zero below $\tau$ and one otherwise;
\emph{boost} reports truth below $\tau$ and one otherwise.
A zero report may still win a zero-bid tie. For each focal user, opponents
use boost on $S\setminus\{f\}$ and truth elsewhere. We compare focal truth,
boost and sandbag from fresh, identical initial states. The timing premium is
$P_\lambda(\tau)=U_f(\mathrm{sandbag})-U_f(\mathrm{boost})$.
We also estimate paired changes $P_\lambda(\tau)-P_{0.5}(\tau)$.
This separates the sign of timing's net benefit from its change with $\lambda$.

\subsection{Non-monetary solutions}
The five history/service extensions are Score-Cap, Karma, Karma-Cap, Rank
and Rank-Cap. Uncapped Karma and Rank disable the service override in the
corresponding capped implementation. Credit-Vickrey is a separate budget extension.
\paragraph{M5c Credit-Vickrey.}
Each user receives income $I=r(k/n)$ before every round. Balances start at
the ceiling $C=20I$. Available credits are $b_i^+=\min(b_i+I,C)$;
the effective bid is $\min(\hat v_i,b_i^+)$. The highest $k$ bids win and pay
the $(k+1)$-st effective bid in credits. Hence
$b_i'=b_i^+-x_i\,\mathrm{price}$. Payments are destroyed, losers pay zero,
and balances cannot become negative. Excess income is discarded; existing
credits do not expire by age. Credits have no direct utility, redemption,
borrowing or transfers. Their cost is reduced future bidding power.
We test $r\in\{0.25,0.5,1,2\}$ rather than assume monetary truthfulness survives.

\paragraph{Service override, Karma-Cap and Rank-Cap.}
For threshold $W$, users with $q_i\ge W$ receive oldest-first priority for
up to $k$ slots. The remaining slots follow the underlying rule.
Starting with zero waits, this guarantees $q_i\le W+\lceil n/k\rceil-1$,
including terminal waits. Once eligible, a user can be preceded by at most
$n-1$ users; each missed round serves $k$ of them, and a served user cannot
overtake it again. The threshold $W$ alone is not a guaranteed wait bound.

Score-Cap applies this override to M4. Karma-Cap starts each balance at two
with ceiling six and caps bids by balances. If $r$ slots remain after the
override, auction winners pay the $(r+1)$-st highest bid among non-forced users.
Forced winners pay the overall $(k+1)$-st bid, limited by their balance.
All credit payments are redistributed equally, then
balances are clipped. This is different from published Karma's credits for
donated capacity \cite{vuppalapati2023karma}; its guarantees do not transfer.
Rank-Cap replaces the raw report by its empirical quantile within that user's
past reports, then applies M4's penalty and the override. Its quantile is
$(N_{<}+N_{=}/2+1/2)/(m+1)$, where $m$ counts past reports and
$N_{<},N_{=}$ count those below or equal to the current report.
Fixed strictly increasing distortions of the
entire report history preserve ranks. Selective reporting and clipping are
outside that invariance property; it is not a truthfulness theorem.
The uncapped rules retain their auction or quantile component but have no
certified finite waiting bound. E9/E10 vary balances, rank blending and wait
thresholds; E12 compares these components with and without the override.

\subsection{Experimental protocol and validation}
The core setting is $n=50$, $k=10$, $T=1000$, independent Uniform$(0,1)$
values and $\rho=0.25$ where strategic. Mixed demand uses equal-sized
Beta$(2,5)$ and Beta$(5,2)$ groups. A Gaussian copula control preserves the
value marginals under persistence; the proposal's AR recursion also shrinks variance.
Each comparison pairs valuation and tie streams. Results use 30 seeds and
10,000 seed-bootstrap resamples. E19/E20 reuse master seeds and are exploratory.
E12--E15 use a separate 30-seed evaluation set, reused descriptively in E16/E17.
E18 uses another 30 evaluation seeds. Adaptive best responses use five disjoint
training traces, eight fixed policies and truthful or doubled-report starts.
Held-out regret maximises seed-averaged unilateral gain over users and policies;
its band is simultaneous within that library. Other intervals are pointwise.
E0 exhausts small discrete cases, including capacity, ties, M4 at $\lambda=0$
and one-round Vickrey deviations. Stateful rules start fresh in each counterfactual.



\newpage
\section{Results}
\subsection{Truthful allocation and sensitivity to exaggeration}
Table~1 compares the five baselines at $n=50$, $k=10$, $T=1000$ under
truthful reports. Values are means across 30 master seeds. Waiting is measured
before each round; $Q_{\max}$ is the mean of the per-episode maxima.
Mixed-demand columns use the two Beta groups, while the other columns use
homogeneous Uniform values. M3 and M5 have the same truthful allocation;
their utilities differ because only M5 charges money.

\begin{center}\small\textbf{Table 1. Truthful welfare, fairness and waiting (E1/E4).}\end{center}
\resulttable{truth}

At $\lambda=1$, Score retains truthful welfare ratio \TruthScoreWR{} under
homogeneous demand. Its starvation rate is \TruthScoreSR{}\%, compared
with \TruthGreedySR{}\% for Greedy. Thus the waiting improvement in this
setting is modest, even though the welfare cost is small. Round-robin supplies
regular service but gives up current-value information. Random allocation
has fair long-run counts without a finite worst-case wait guarantee.

Heterogeneity changes the interpretation of fairness. In the mixed population,
Greedy's normalised-benefit index is \MixedGreedyJB{}, compared with
\MixedScoreJB{} for Score. This improvement has a larger welfare cost:
Score's mixed-demand welfare ratio is \MixedScoreWR{}. The claim that a
history correction is nearly free therefore applies to the homogeneous
truthful setting, not to every workload.

\resultfigure{m4_inflation}{\textbf{Figure 1.} M4 unilateral gains under capped exaggeration (E7).
Opponents use the same exaggeration policy within each cell; $\rho=0.25$.
Bars are pointwise 95\% seed-bootstrap intervals.}

Figure~1 shows why the penalty needed against one inflation rule cannot be
treated as a general deterrence threshold. At $\lambda=2$, doubling values
has mean gain \DoubleTwo{}, while the milder factor $c=1.5$ has gain
\MediumTwo{}. At $\lambda=5$, $c=1.5$ has gain \MediumFive{}, but
$c=1.25$ still yields \MildFive{}. These are policy-specific mean gains,
not claims that every user benefits or that all other deviations are deterred.
The figure also changes the opponents' policy when $c$ changes, so it does
not rank the policies as best responses to one common opponent population.

\paragraph{Other workload and implementation checks.}
Scarcity/frontier and strategic-fraction sweeps (E1--E3) compare welfare,
fairness and policy-specific gain separately. E5 isolates persistence from
shrinking variance in the proposal recursion: with marginal-preserving
persistence at $\alpha=0.9$, Score's starvation reaches \PersistentScoreSR{}\%.
Equal long-run counts need not give regular service. E6 benchmarks $n=10$--500.
Saved Score runtime at $n=500$, $k=100$ averages \ScoreRuntime{} microseconds
per round (30 seeds, best of three after warm-up), excluding value generation.
Peak Python heap includes the seed package and episode records, sampled on
three seeds. These machine-dependent simulator benchmarks do not measure a production scheduler.

\newpage
\subsection{Core timing experiment: does suppressing low reports help?}
E19 runs at the full core population size, rather than relying on the small
rollout diagnostic. Table~2 distinguishes two quantities at $\tau=0.75$.
The original population-profile gain compares sandbag to truth while other
strategic users also sandbag. The controlled gain holds those opponents at
boost. Both are valid unilateral comparisons, but they describe different games.
Subtracting gains from the two opponent populations would not isolate timing.

\begin{center}\small\textbf{Table 2. M4 timing at $\tau=0.75$ (E19); 95\% intervals.}\end{center}
\resulttable{timing}

The last column is the controlled premium $P_\lambda(0.75)$, which changes
only the focal policy from boost to sandbag. Its value is \TimingLow{} at
$\lambda=0.5$ and \TimingHigh{} at $\lambda=5$. The seed-paired increase
between those penalties is \TimingTrend{}. These intervals support the
direction of hypothesis H2 for this threshold and comparison.
They do not establish monotonicity for all penalties or reporting rules.

\resultfigure{m4_timing}{\textbf{Figure 2.} Controlled timing premiums across all tested thresholds.
Each arm has the same opponents, values, ties and initial state. The high-threshold
policy is shown on a separate y scale so its losses remain visible.}

Threshold choice matters. At $\tau=0.5$, suppressing low reports has little
effect because those rounds rarely provide useful service in these settings.
At $\tau=0.9$, it gives up sufficiently valuable rounds that its controlled
premium remains strongly negative throughout the grid. Thus a stronger
history penalty does not make every selective-reporting rule profitable.
At $\tau=0.75$, the benefit of suppression is small relative to total
episode utility; much of the policy's gain is already available to boost.

The timing comparison includes Greedy and monetary Vickrey as controls.
Their controlled premiums are non-positive: suppressing low reports has no
future allocation-history benefit in these rules under the fixed,
value-dependent opponent policies. Credit-Vickrey at rate $0.5$ instead
has premium \CreditTiming{}. There, changing current service also changes
future bidding resources. This is a separate dynamic budget effect.

\paragraph{What H2 means here.}
The experiment measures the net effect of suppressing reports: current
benefit forgone plus all downstream allocation effects. It is not an additive
causal decomposition of the full policy gain into independent mechanisms.
The displayed comparison was examined on existing master seeds and the
threshold grid is exploratory. A positive interval is evidence for a tested
deviation on this workload distribution, not an unrestricted optimal policy.
The code saves each focal user's three utilities before averaging, and the
audit checks these identities and the paired premium changes directly.

\newpage
\subsection{Vickrey with renewable credits}
E20 replaces utility-costly payments with the budget model in Section~3.
Table~3 uses $n=50$, $k=10$, $T=1000$ and the same 30 master seeds.
The fairness and waiting columns are from truthful runs. The final columns
are unilateral gains in the reporting policy's own strategic population,
with nominal fraction $\rho=0.25$. They are not direct policy rankings under
one shared opponent profile.

\begin{center}\small\textbf{Table 3. Credit income, truthful outcomes and strategic gains (E20).}\end{center}
\resulttable{credits}

A tight quota has a substantial efficiency cost. At rate $0.25$, the truthful
welfare ratio is \CreditTightWR{}, while the mean per-episode maximum wait
is \CreditTightQ{} rounds, compared with \TruthGreedyQ{} for truthful
Greedy. The allocation-count Jain index is close to one. This is observed
equality in the sampled workloads, not a guaranteed bound on individual waits.
At rate $0.5$, truthful welfare rises to \CreditHalfWR{}.

\resultfigure{credit_budget}{\textbf{Figure 3.} Credit-Vickrey's efficiency--fairness trade-off
and policy-specific unilateral gains. Bars are pointwise 95\% intervals.}

The tested inflation policies have negative mean gains at the tighter rates,
but sandbag at $\tau=0.75$ remains profitable. With looser budgets, some
inflation policies also have positive mean gains. No tested income rate
eliminates every profitable tested deviation. The controlled timing arm in
E19 supplies a fixed-opponent check at rate $0.5$; E20's different policy
populations alone cannot identify how much of each gain is caused by saving credits.

\paragraph{Why monetary truthfulness does not carry over.}
Monetary Vickrey subtracts the price directly from utility. M5c instead
subtracts credits from a future bidding budget. With a sufficiently loose
finite budget, M5c has Greedy's allocation but zero utility payments; then
inflating a report can obtain a positive-value GPU that truth would lose.
The test suite contains this one-round counterexample. Under a tight budget,
future opportunities matter, so neither truthful bidding nor always bidding
the maximum is justified by the one-round monetary proof.

\paragraph{Scope of the extension.}
The experiment evaluates one explicit budget design. The initial balance,
income, ceiling and destruction of payments are part of the model, not
features shared by all cluster-credit systems. We do not infer results for
transferable credits, redistribution, borrowing or different expiry rules.
The budget has no service override and supplies no hard waiting guarantee.

\newpage
\subsection{Five history/service extensions and stronger attacks}
Table~4 includes all five extensions and uncapped Score as their baseline.
Score and Rank variants use $\lambda=1$. Capped rules use $W=10$ and share the certified wait bound
$10+\lceil50/10\rceil-1=14$. Truthful columns use the named workload;
the final column uses homogeneous demand, $c=1.25$ and $\rho=0.25$.
Persistent values use the marginal-preserving copula with $\alpha=0.9$.
These results use the fresh extension seeds, not the core master seeds.

\begin{center}\small\textbf{Table 4. Capped and uncapped component outcomes (E12).}\end{center}
\resulttable{extensions}

Removing the override raises welfare in each listed workload. Score and Rank
variants retain positive mild-inflation gains; both Karma variants have negative
means in this population. Karma-Cap reverses when all opponents inflate:
at $c=2$, $\rho=1$, its gain is \KarmaPopulationReversal{}.
Equal caps therefore do not ensure incentive compatibility. Heterogeneous
and persistent demand increase the welfare cost.

\paragraph{Corrected planning attack (E15).}
A focal user searches $\{0,0.1,\ldots,1\}$ plus exact truth. Candidates share
sampled futures, truncated at episode end; hypothetical counts, balances and
rank histories advance. Only the focal current value is observed. Opponent
values are sampled and focal continuation is truthful. Figure~4 uses $n=10$,
$k=2$, $T=120$, horizon 40, 80 futures per candidate, truthful opponents and
30 extension seeds.

\resultfigure{planning_attack}{\textbf{Figure 4.} Corrected rollout gains (E15).
Gain is measured over 120 rounds, unlike the 1,000-round gains above.
Negative or zero-crossing intervals do not prove strategyproofness.}

The planner gains \ScorePlanner{} against Score at $\lambda=2$ and
\RankPlanner{} against Rank-Cap. Thus the tested fixed inflation results
do not exhaust these mechanisms' vulnerabilities. Karma-Cap's interval
includes zero in this setting; this is a limited attack outcome, not a guarantee.
Historical E3b/E3c/E11 checkpoints used different attacker assumptions;
their numerical claims are superseded here by corrected E15.

\paragraph{Welfare relative to the same service constraint (E14).}
For $B=W+\lceil n/k\rceil-1$, an offline binary program requires each user
to receive service in every $B+1$-round window, including the last window.
The oracle knows future values. All \CertifiedSolves{} unique small-instance
solves across the full grid are certified optimal. Table~5 shows the mixed
case at $n=6$, $k=2$, $T=40$, $W=6$ and $B=8$.

\begin{center}\small\textbf{Table 5. Welfare under an identical service requirement (E14).}\end{center}
\resulttable{oracle}

Constrained efficiency is $W_{\rm online}/W_B^*$. Service loss is
$(W^*-W_B^*)/W^*$; remaining loss is $(W_B^*-W_{\rm online})/W^*$.
The latter includes the online information disadvantage. A service guarantee
therefore does not explain the entire gap from the future-aware optimum.

\newpage
\subsection{Adaptive policies and failure diagnostics}
E13 trains asynchronous best responses within an eight-policy library.
E18 extends training from four to twenty sweeps at $n=6$, $k=2$, $T=120$;
Table~6 evaluates the frozen profiles on 30 independent follow-up seeds.
Karma-Cap converges on the five training traces at sweep 18 from both starts.
Score-Cap and Rank-Cap reach the twenty-sweep limit. Held-out regret remains
uncertain; training convergence establishes neither population nor unrestricted
equilibrium. The library excludes E19's value-threshold policies.

\begin{center}\small\textbf{Table 6. Restricted held-out regret; simultaneous 95\% bands (E18).}\end{center}
\resulttable{adaptive}

E16 replays Rank-Cap's fixed attack reports using a private truthful reference
for the focal user's rank history.
Updating the history adds \RankLedgerEffect{} utility relative to that replay.
This privileged diagnostic is not a deployable rule; the attack also profits
with the truthful reference. E17 decomposes Karma-Cap's $c=2$, $\rho=1$
gain into forced-service \KarmaForcedGain{} and auction \KarmaAuctionGain{}
benefits. Uncapped Karma also reverses sign: \KarmaUncappedReversal{}.
The override therefore does not alone explain the reversal. These reused-seed
diagnostics describe accounting and history effects, not unique causal or
equilibrium explanations.

\section{Discussion and limitations}
\paragraph{Returning to the proposal's hypotheses.}
\textbf{H1:} Moderate history penalties cost little homogeneous truthful welfare,
but reduce starvation only modestly. \textbf{H2:} Timing's premium increases
between tested penalties at $\tau=0.75$; the costly $\tau=0.9$ policy prevents
generalisation. \textbf{H3:} Round-robin trades welfare for equal counts and
also leads mean normalised-benefit fairness in the mixed case.
\textbf{H4:} Monetary Vickrey supplies the one-round theory control; its proof
does not cover credits. M3 and M4 admit profitable tested deviations.

\paragraph{Limits of the evidence.}
Identical GPUs, one-round service and additive synthetic values omit arrivals,
placement and multi-GPU jobs; welfare is not training throughput. Fairness
indices capture specific equality objectives. Rollouts use marginal futures
and truthful continuation, not optimal best responses. E19/E20 reuse master
seeds; their exploratory pointwise intervals do not cover all inspected grid
comparisons. None of the tested non-monetary rules is established as dynamically
strategyproof.

\section{Contribution}
All four primary assignments were fulfilled.
\begin{center}\small
\begin{tabular}{p{1.05in}p{4.05in}p{0.65in}}
\toprule Member & Assigned task and delivered component & Status \\
\midrule
Aayush Kuloor & Valuations, histories, paired random streams, configuration,
result storage and end-to-end runner; seed and reproduction support. & Fulfilled \\
Harsh Dhru & Mechanisms, allocation/payment interfaces, seeded ties and exhaustive
E0 checks of invariants and the Vickrey benchmark. & Fulfilled \\
Raj Modi & Bounded reporting policies, rollout diagnostic, welfare/fairness/waiting
and strategic metrics, with hand-computed validation. & Fulfilled \\
Kalp Shah & Parameter sweeps, paired bootstrap intervals, Pareto analysis,
figures, result tables and sensitivity/failure documentation. & Fulfilled \\
\bottomrule
\end{tabular}\end{center}

\section{Repository and reproducibility}
\repo

Install \texttt{requirements-lock.txt} and run \verb|python run_all.py|
to rerun experiments and regenerate result plots and tables. Configurations
and seed lists are included. Run \verb|python -m pytest| for validation.
\verb|python run_all.py --quick| is a smoke test, not the reported evaluation.
\label{lastbodypage}

\clearpage
\bibliographystyle{unsrt}
\bibliography{final_references}
\end{document}
